% Generated draft: requires mathematical and source review.
\documentclass[11pt]{article}
\usepackage[margin=25mm]{geometry}
\usepackage{fontspec}
\setmainfont{Latin Modern Roman}
\usepackage{amsmath,amssymb,mathtools,unicode-math}
\setmathfont{Latin Modern Math}
\usepackage{enumitem,needspace,xurl,hyperref}
\setlength{\parindent}{0pt}
\setlength{\parskip}{5pt}
\setlength{\emergencystretch}{3em}
\newcommand{\N}{\mathbb{N}}
\newcommand{\Z}{\mathbb{Z}}
\newcommand{\Q}{\mathbb{Q}}
\newcommand{\R}{\mathbb{R}}
\newcommand{\C}{\mathbb{C}}
\newcommand{\F}{\mathbb{F}}
\newcommand{\A}{\mathbb{A}}
\newcommand{\PP}{\mathbb{P}}
\newcommand{\E}{\mathbb{E}}
\newcommand{\eps}{\varepsilon}
\newcommand{\dd}{\,\mathrm{d}}
\begin{document}
\section*{UnsolvedMath ID: 30003001}
\textbf{Source problem code:} catalogue entry
\subsection{Definitions and mathematical statement}
A Dirac peak means a point mass, conventionally written \(m\delta_a\), where \(a\) is its location and \(m\) its mass. The Dirac distribution on \(\R\) is characterized by
\[
\langle\delta_a,\varphi\rangle=\varphi(a),\qquad \varphi\in C_c^\infty(\R).
\]
This notation does not specify the spatial domain, admissible masses, or dependent variables of the cytoskeleton model.

The supplied problem asks: for the spatially one-dimensional hyperbolic--parabolic cytoskeleton model referred to in the cited report, can the explicitly constructed Dirac peaks, moving profiles, and shocks arise dynamically from a suitable class of smooth initial data?

The intended model and target solutions are defined only by reference to that report. The supplied record contains no evolution equations, coefficients, spatial domain, boundary conditions, admissibility conditions, or formulas for the constructed peaks, profiles, and shocks. It also does not specify the solution concept after singularity formation or whether ``arise dynamically'' means finite-time formation, asymptotic convergence, or another relation. No convergence topology or quantification over the target families is provided. Consequently, the question above preserves the supplied target, but a fully quantified mathematical formulation cannot be recovered from this record alone.
\subsection{Short English statement}
Can suitable smooth initial data for the referenced one-dimensional hyperbolic--parabolic cytoskeleton model dynamically produce its explicitly constructed Dirac peaks, moving profiles, and shocks? The governing equations, target families, and precise meaning of dynamical formation must be recovered from the cited report before this becomes a fully specified mathematical problem.
\subsection{Sources}
\begin{itemize}
\item Local UnsolvedMath registry, supplied source record: permanent ID \(30003001\), title \emph{Formation of Singular Cytoskeleton Profiles from Smooth Data}, source problem code \texttt{catalogue entry}, collection \emph{Oberwolfach Reports Open Problems}. The record gives proposed year 2015 and names no individual proposer.
\item Supplied report citation: \emph{Geometric Partial Differential Equations: Surface and Bulk Processes} (2016), Oberwolfach Reports, pp.\ 3101--3178, DOI \texttt{10.4171/owr/2015/55}. Supplied source URL and literature-source link, labeled ``Geometric PDE Oberwolfach report'': \url{https://doi.org/10.4171/owr/2015/55}. The supplied background describes these reports as records of mathematical workshops and their open-problem discussions.
\item Supplied literature assessment, checked \textbf{2026-08-21}: the record labels the problem open and states that the report explicitly asks whether the displayed shocks, peaks, or profiles can arise dynamically from suitable smooth initial data; its curation identified no definitive general answer in the literature. This is supplied dated curation, not an independently verified literature assessment. The supplied background expressly requires independent verification of time-sensitive status and cited claims.
\end{itemize}
\end{document}
